\section{Conclusão}
\label{sec_conclusao}

Este trabalho apresentou o projeto e a implementação de um serviço de metadados compatível com cloud-init para o Incus. O serviço preenche uma lacuna no ecossistema de gerenciamento de contêineres e máquinas virtuais, oferecendo um endpoint de metadados independente que as instâncias podem consultar para obter suas configurações de inicialização, de forma análoga ao que provedores de nuvem pública disponibilizam nativamente. Foram implementados os endpoints REST exigidos pelas fontes de dados do cloud-init (\texttt{meta-data}, \texttt{user-data}, \texttt{vendor-data} e \texttt{network-config}), além de um sistema de sincronização orientado a eventos que mantém os dados das instâncias atualizados no banco de dados. A identificação das instâncias por endereço IP dispensa agentes dentro das instâncias; a única configuração necessária é o parâmetro \texttt{seedfrom} do \textit{datasource} NoCloud, definido uma única vez na imagem-base. O serviço implementa ainda suporte a consenso distribuído via RAFT, habilitando implantações de alta disponibilidade com múltiplos nós coordenados por eleição de líder, replicação de log e aplicação determinística de comandos de escrita por meio de uma máquina de estados finita (FSM).

A avaliação experimental verificou a corretude funcional e o desempenho do serviço em cenários com diferentes números de instâncias simultâneas (de 5 a 200 contêineres), bem como o mecanismo de alta disponibilidade baseado em consenso RAFT em um \textit{cluster} de três nós. Os resultados validaram a viabilidade da abordagem em ambientes Incus reais: as inicializações de \textit{cloud-init} completaram-se com sucesso, as latências de resposta situaram-se na faixa de dezenas de milissegundos e a reeleição de líder após falha ocorreu em cerca de dois segundos. Algumas limitações foram identificadas. O uso do SQLite impõe restrições à concorrência de escrita em cenários de alta carga, ainda que não se tenha observado degradação até 200 instâncias. O conjunto de campos de metadados suportados também representa um subconjunto do que plataformas como AWS e GCP oferecem.

Como trabalhos futuros, destacam-se as seguintes direções: (i) eliminar a dependência de uma imagem-base pré-configurada, por exemplo por meio de um \textit{datasource} nativo do Incus no cloud-init ou da passagem do parâmetro \texttt{seedfrom} pela linha de comando do \textit{kernel} em máquinas virtuais; (ii) ampliar o suporte a módulos e campos do cloud-init para maior compatibilidade com cargas de trabalho existentes; (iii) implementar autenticação baseada em sessões inspirada no IMDSv2~\citep{aws_imdsv2} para mitigar riscos de SSRF~\citep{mitre_t1552_005}; (iv) substituir o mecanismo de sincronização por \textit{polling} pela API de eventos do Incus, obtendo sincronização em tempo real com menor sobrecarga; (v) realizar \textit{benchmarks} de desempenho em cenários de maior escala, com repetições suficientes para análise estatística; e (vi) validar a aplicação da configuração de rede entregue pelo serviço em imagens com a gestão de rede do cloud-init habilitada. Estas melhorias consolidariam o serviço como uma solução robusta para ambientes privados que demandam automação de infraestrutura baseada em cloud-init.
